% ============================================================
%  SY02-36 约化群的 Kloosterman 层 —— 导言区
%  原文：J. Heinloth, B.-C. Ngô, Z. Zhiwei Yun,
%        "Kloosterman sheaves for reductive groups"
%        arXiv:1005.2765v1 (16 May 2010), 58 页
% ============================================================

\usepackage[a4paper,margin=2.6cm]{geometry}

\PassOptionsToPackage{leqno}{amsmath}% 原文公式号在左
\usepackage{amsmath,amssymb,amsfonts,amsthm}
\usepackage{mathtools}
\usepackage{mathrsfs}
\usepackage{bm}
\usepackage{graphicx}
\usepackage{booktabs}
\usepackage{array}
\usepackage{tabularx}
\usepackage{enumitem}
\usepackage{xcolor}
\usepackage[most]{tcolorbox}
\usepackage{xspace}
\usepackage{tikz}
\usetikzlibrary{cd}
\usepackage{tikz-cd}
\usepackage{microtype}

\usepackage[bookmarks=true,bookmarksnumbered=true,
  colorlinks=true,linkcolor=blue!45!black,citecolor=blue!45!black,
  urlcolor=teal!70!black,filecolor=teal!70!black,
  pdftitle={约化群的 Kloosterman 层},
  pdfauthor={J. Heinloth, B.-C. Ngô, Z. Yun（译者：数译）}]{hyperref}

\linespread{1.22}
\setlength{\parskip}{.25em}
\setlength{\emergencystretch}{3em}
\sloppy

% ---------- 中文排版 ----------
\xeCJKsetup{CJKecglue={\hskip 0.24em plus 0.05em minus 0.05em}}
\ctexset{
  section    = {name={,}, number=\thesection,    aftername={\enspace}},
  subsection = {name={,}, number=\thesubsection, aftername={\enspace}},
  subsubsection = {name={,}, number=\thesubsubsection, aftername={\enspace}},
}
% 引言为「第 0 节」：section 计数器从 -1 起，显示为 0
\renewcommand{\thesection}{%
  \ifnum\value{section}<0 0\else\arabic{section}\fi}
\renewcommand{\thesubsection}{\thesection.\arabic{subsection}}
\renewcommand{\thesubsubsection}{\thesubsection.\arabic{subsubsection}}

% ---------- 编号 ----------
% 原文：主定理 Theorem 1–4 全局编号（用 mainthm 手写号）；
% 其余定理类与编号公式共用一个按节递增的计数器。
\numberwithin{equation}{section}
\theoremstyle{plain}
\newtheorem{theorem}[equation]{定理}
\newtheorem{lemma}[equation]{引理}
\newtheorem{proposition}[equation]{命题}
\newtheorem{corollary}[equation]{推论}
\newtheorem{definition}[equation]{定义}
\newtheorem{example}[equation]{例}
\newtheorem{remark}[equation]{注}
\newtheorem{claim}[equation]{断言}
\newtheorem{notation}[equation]{记号}
\newtheorem{conjecture}[equation]{猜想}
\newtheorem{fact}[equation]{事实}
\theoremstyle{definition}

% 主定理：全局号 1–4，不占节内计数器
\newenvironment{mainthm}[1][]{%
  \par\addvspace{\medskipamount}%
  \noindent\textbf{定理 #1.}\ \itshape\ignorespaces
}{%
  \par\addvspace{\medskipamount}%
}
\newenvironment{mainthmstar}[1][]{%
  \par\addvspace{\medskipamount}%
  \noindent\textbf{定理 #1}\ \itshape\ignorespaces
}{%
  \par\addvspace{\medskipamount}%
}

\setcounter{tocdepth}{2}
\setcounter{secnumdepth}{3}

% ---------- 译注 ----------
\newcounter{transnote}
\newtcolorbox{yizhu}{
  colback=orange!4, colframe=orange!55!black, boxrule=.4pt,
  left=6pt, right=6pt, top=4pt, bottom=4pt,
  fontupper=\small,
  before skip=6pt, after skip=6pt,
}
\makeatletter
\renewenvironment{yizhu}{%
  \refstepcounter{transnote}%
  \begin{tcolorbox}[colback=orange!4, colframe=orange!55!black, boxrule=.4pt,
    left=6pt, right=6pt, top=4pt, bottom=4pt, fontupper=\small,
    before skip=6pt, after skip=6pt]%
  \noindent\textbf{译注 \thetransnote：}\ignorespaces
}{%
  \end{tcolorbox}%
}
\makeatother

% 行内译注：橙色小字脚注
\newcommand{\zhnote}[1]{%
  \footnote{\textcolor{orange!70!black}{【译注】}#1}}

% 路线注：每块开头标原文页码
\newcommand{\routenote}[1]{%
  \par\noindent{\footnotesize\color{gray}［对应原文 #1］}\par\addvspace{.2em}}

% ---------- 常用记号 ----------
\newcommand{\Bun}{\mathrm{Bun}}
\newcommand{\Gr}{\mathrm{Gr}}
\newcommand{\Loc}{\mathrm{Loc}}
\newcommand{\Rep}{\mathrm{Rep}}
\renewcommand{\Vec}{\mathrm{Vec}}
\newcommand{\Perv}{\mathrm{Perv}}
\newcommand{\IC}{\mathrm{IC}}
\newcommand{\AS}{\mathrm{AS}}
\newcommand{\Kum}{\mathrm{Kum}}
\newcommand{\Swan}{\mathrm{Swan}}
\newcommand{\Kl}{\mathrm{Kl}}
\newcommand{\Hk}{\mathrm{Hk}}
\newcommand{\Hecke}{\mathcal{H}\mathrm{ecke}}
\newcommand{\Fl}{\mathrm{Fl}}
\newcommand{\Aut}{\mathrm{Aut}}
\newcommand{\Out}{\mathrm{Out}}
\newcommand{\Stab}{\mathrm{Stab}}
\newcommand{\Sat}{\mathcal{S}\mathrm{at}}
\newcommand{\semiSat}{\mathcal{S}}
\newcommand{\cO}{\mathcal{O}}
\newcommand{\cG}{\mathcal{G}}
\newcommand{\cS}{\mathcal{S}}
\newcommand{\cC}{\mathcal{C}}
\newcommand{\cT}{\mathcal{T}}
\newcommand{\cF}{\mathcal{F}}
\newcommand{\cL}{\mathcal{L}}
\newcommand{\cE}{\mathcal{E}}
\newcommand{\cH}{\mathcal{H}}
\newcommand{\ck}{\check}
\newcommand{\ckG}{\check G}
\newcommand{\ckT}{\check T}
\newcommand{\ckg}{\check{\mathfrak g}}
\newcommand{\ckb}{\check{\mathfrak b}}
\newcommand{\ckn}{\check{\mathfrak n}}
\newcommand{\LG}{{}^L G}
\newcommand{\frakg}{\mathfrak{g}}
\newcommand{\frakt}{\mathfrak{t}}
\newcommand{\frakb}{\mathfrak{b}}
\newcommand{\frakn}{\mathfrak{n}}
\newcommand{\frakh}{\mathfrak{h}}
\newcommand{\fraku}{\mathfrak{u}}
\newcommand{\frakm}{\mathfrak{m}}
\newcommand{\frakz}{\mathfrak{z}}
\newcommand{\ev}{\mathrm{ev}}
\newcommand{\pr}{\mathrm{pr}}
\newcommand{\ad}{\mathrm{ad}}
\newcommand{\Ad}{\mathrm{Ad}}
\newcommand{\mult}{\mathrm{mult}}
\newcommand{\inv}{\mathrm{inv}}
\newcommand{\St}{\mathrm{Std}}
\newcommand{\dd}{\mathrm{d}}
\newcommand{\sgn}{\mathrm{sgn}}
\newcommand{\CC}{\mathbb{C}}
\newcommand{\PP}{\mathbb{P}}
\newcommand{\GG}{\mathbb{G}}
\newcommand{\Aff}{\mathbb{A}}
\newcommand{\ZZ}{\mathbb{Z}}
\newcommand{\QQ}{\mathbb{Q}}
\newcommand{\FF}{\mathbb{F}}
\newcommand{\Ql}{\bar{\mathbb{Q}}_\ell}
\newcommand{\Wt}{\widetilde{W}}
\newcommand{\It}{I^t}
\newcommand{\Ip}{I^+}
\newcommand{\Satgeo}{\Sat^{\mathrm{geo}}}
\newcommand{\Grtriv}{\Gr^{\mathrm{triv}}}
\newcommand{\KlG}{\Kl_{\ckG}}
\newcommand{\Klg}{\Kl_{\mathfrak g}}
\newcommand{\red}{\mathrm{red}}
\newcommand{\der}{\mathrm{der}}
\newcommand{\loc}{\mathrm{loc}}
\newcommand{\rs}{r_s}
\newcommand{\rl}{r_\ell}

\pagestyle{plain}
